\title{Lecture 8: 2D Transformation with Matrix}
\frame{\titlepage}

%%%%%%%%%%%%%
\section[Intro. to Matrix]{Introduction to Matrix}
\frame{\sectionpage}
%%%%%%%%%%%%%

\begin{frame}
\frametitle{Matrix Definition}
\begin{defn}{}{}
A matrix is an ordered collection of numbers. A 2 by 2 matrix is of the form
\[ \begin{pmatrix} a & b \\ c & d \end{pmatrix}\]
where $a,b,c,d$ are real numbers.
\end{defn}
In general, we can define $m$ by $n$ matrix in the same way, but we only need to focus on a 2 by 2 matrix for today's topic.
\end{frame}


\begin{frame}
\frametitle{Matrix Operation}
Given two matrices $A = \begin{pmatrix} a & b \\ c & d \end{pmatrix}$ and $A' = \begin{pmatrix} a' & b' \\ c' & d' \end{pmatrix}$, a vector $\vec{v} = \cvec{v_1}{v_2}$, and a real number $k$. We can define the following operations.
\begin{itemize}
\item Scalar multiplication: $kA$
\item Matrix addition: $A + A'$
\item Matrix-vector multiplication: $A\vec{v}$
\item Matrix multiplication: $AA'$
\end{itemize}
\end{frame}


\begin{frame}
\frametitle{Scalar Multiplication}
\begin{defn}{Scalar multiplication}{}
\begin{equation}
\red{k}\begin{pmatrix} a & b \\ c & d \end{pmatrix}= \begin{pmatrix} \red{k}a & \red{k}b \\ \red{k}c & \red{k}d \end{pmatrix}
\end{equation}
\end{defn}
\end{frame}


\begin{frame}[t]
\begin{ex}
Compute the followings.
\begin{tasks}(2)
\task $4\begin{pmatrix} 1 & 0 \\ -2 & 3 \end{pmatrix}$
\task $-2\begin{pmatrix} 3 & -1 \\ 0 & 3 \end{pmatrix}$
\end{tasks}
\end{ex}
\begin{sol}
\begin{tasks}(1)
\task $\red{4}\begin{pmatrix} 1 & 0 \\ -2 & 3 \end{pmatrix} = \begin{pmatrix} \red{4} \times 1 & \red{4} \times 0 \\ \red{4} \times -2 & \red{4} \times 3 \end{pmatrix} = \begin{pmatrix} 4 & 0 \\ -8 & 12 \end{pmatrix}$
\task $\red{-2}\begin{pmatrix} 3 & -1 \\ 0 & 3 \end{pmatrix} = \begin{pmatrix} \red{(-2)} \times 3 & \red{(-2)} \times (-1) \\ \red{(-2)} \times 0 & \red{(-2)} \times 3 \end{pmatrix} = \begin{pmatrix} -6 & 2 \\ 0 & -6 \end{pmatrix}$
\end{tasks}
\end{sol}
\end{frame}

\begin{frame}
\frametitle{Matrix Addition}
Matrix addition is done entry by entry.
\begin{defn}{Matrix addition}{}
\begin{equation}
\begin{pmatrix} a & b \\ c & d \end{pmatrix} +  \begin{pmatrix} a' & b' \\ c' & d' \end{pmatrix} = \begin{pmatrix} a+a' & b + b' \\ c+ c' & d + d' \end{pmatrix}
\end{equation}
\end{defn}
\end{frame}


\begin{frame}[t]
\begin{ex}
Compute the followings.
\begin{tasks}(2)
\task $\begin{pmatrix} 0 & 1 \\ 2 & 3 \end{pmatrix} +  \begin{pmatrix} 7 & -6 \\ -5 & 4 \end{pmatrix}$
\task $\begin{pmatrix} 0 & 1 \\ 2 & 3 \end{pmatrix} +  \begin{pmatrix} 7 & -6 \\ -5 & 4 \end{pmatrix}$
\end{tasks}
\end{ex}
\begin{sol}
\begin{tasks}(1)
\task $\begin{pmatrix} 0 & 1 \\ 2 & 3 \end{pmatrix} +  \begin{pmatrix} 7 & -6 \\ -5 & 4 \end{pmatrix} = \begin{pmatrix} 0+7 & 1+(-6) \\ 2+(-5) & 3+4 \end{pmatrix} = \begin{pmatrix} 7 & -5 \\-3 & 7 \end{pmatrix}$
\task $\begin{pmatrix} 0 & 1 \\ 2 & 3 \end{pmatrix} +  \begin{pmatrix} 7 & -6 \\ -5 & 4 \end{pmatrix} = \begin{pmatrix} 0+7 & 1+(-6) \\ 2+(-5) & 3+4 \end{pmatrix} = \begin{pmatrix} 7 & -5 \\-3 & 7 \end{pmatrix}$
\end{tasks}
\end{sol}
\end{frame}


\begin{frame}
\frametitle{Matrix-Vector Multiplication}
\begin{defn}{Matrix-Vector Multiplication}{}
\begin{equation}
\begin{pmatrix} \red{a} & \green{b} \\ \red{c} & \green{d} \end{pmatrix}\cvec{v_1}{v_2}  = v_1 \cvec{\red{a}}{\red{c}} + v_2 \cvec{\green{b}}{\green{d}} = \cvec{\red{a}v_1 + \green{b}v_2}{\red{c}v_1 +\green{b}v_2}
\end{equation}
\end{defn}
\end{frame}


\begin{frame}[t]
\begin{ex}
Compute the followings.
\begin{tasks}(2)
\task $\begin{pmatrix} \red{2} & \green{3} \\ \red{-1} & \green{0} \end{pmatrix}\cvec{4}{5}$
\task $\begin{pmatrix} \red{-3} & \green{1} \\ \red{5} & \green{-2} \end{pmatrix}\cvec{1}{0}$
\end{tasks}
\end{ex}
\begin{sol}
\begin{tasks}(1)
\task $\begin{pmatrix} \red{2} & \green{3} \\ \red{-1} & \green{0} \end{pmatrix}\cvec{4}{5} = 4 \cvec{\red{2}}{\red{-1}} + 5 \cvec{\green{3}}{\green{0}}
= \cvec{8}{-4} + \cvec{15}{0} = \cvec{23}{-4}$
\task $\begin{pmatrix} \red{-3} & \green{1} \\ \red{5} & \green{-2} \end{pmatrix}\cvec{1}{0} = 1 \cvec{\red{-3}}{\red{1}} + 0 \cvec{\green{5}}{\green{2}}
= \cvec{-3}{1} + \cvec{0}{0} = \cvec{-3}{1}$
\end{tasks}
\end{sol}
\end{frame}

\begin{frame}
\frametitle{Motivation for Matrix-Vector Multiplication}
Every single math student is perplexed by the formula
\[ \begin{pmatrix} \red{a} & \green{b} \\ \red{c} & \green{d} \end{pmatrix}\cvec{\blue{v_1}}{\orange{v_2}}  = \cvec{\red{a}\blue{v_1} + \green{b}\orange{v_2}}{\red{c}\blue{v_1} +\green{b}\orange{v_2}}.\]
Why is it defined this way? Won't it be much simpler if we define it naturally as
\[ \begin{pmatrix} a & b \\ c & d \end{pmatrix}\cvec{v_1}{v_2} = \begin{pmatrix} av_1 & bv_1 \\ cv_2 & dv_2 \end{pmatrix}\]
or 
\[ \begin{pmatrix} a & b \\ c & d \end{pmatrix}\cvec{v_1}{v_2} =  \begin{pmatrix} av_1 + bv_1 \\ cv_2 + dv_2 \end{pmatrix}?\]
\end{frame}

\begin{frame}
The answer is: this is the \textit{natural} way to define the multiplication. 

Notice what happens when we multiply the matrix by the unit vector $\vec{i} = \cvec{1}{0}$.
\[ \begin{pmatrix} \red{a} & \green{b} \\ \red{c} & \green{d} \end{pmatrix}\cvec{1}{0}  = \cvec{\red{a}\times 1 + \green{b} \times 0}{\red{c} \times 1 +\green{b} \times 0} = \cvec{\red{a}}{\red{c}}.\]
So the matrix transform $\vec{i}$ to the vector $\cvec{\red{a}}{\red{c}}$, which is the first column of the matrix. 
\end{frame}

\begin{frame}
Likewise, for the unit vector $\vec{j} = \cvec{0}{1}$, we have
\[ \begin{pmatrix} \red{a} & \green{b} \\ \red{c} & \green{d} \end{pmatrix}\cvec{0}{1}  = \cvec{\red{a}\times 0 + \green{b} \times 1}{\red{c} \times 0 +\green{b} \times 1} = \cvec{\green{b}}{\green{d}}.\]
So the matrix transform $\vec{j}$ to the vector $\cvec{\green{b}}{\green{d}}$, which is the second column of the matrix.
\end{frame}

\begin{frame}
Combining the equations together and using the fact that the matrix multiplication is linear, we can visually define.
\begin{eq*} 
\begin{pmatrix} \red{a} & \green{b} \\ \red{c} & \green{d} \end{pmatrix}\cvec{\blue{v_1}}{\orange{v_2}}  &=& \begin{pmatrix} \red{a} & \green{b} \\ \red{c} & \green{d} \end{pmatrix} \left( \cvec{\blue{v_1}}{0} + \cvec{0}{\orange{v_2}} \right) \\
&=& \cvec{\red{a}\blue{v_1} + \green{b}\orange{v_2}}{\red{c}\blue{v_1} +\green{b}\orange{v_2}}
\end{eq*}
\end{frame}

\begin{frame}
\addfig{0.9}{2dtransform_multiplication_explained}{The grid transformed by the matrix. \href{https://www.desmos.com/calculator/wt7lapyxi4}{Click for interactive tool at desmos.}}{}
\end{frame}

\begin{frame}
\begin{exercisebox}{}{}
Compute the following matrix operations.
\begin{tasks}(2)
\task $3 \begin{pmatrix} 1 & 2 \\ -1 & 0 \end{pmatrix}$
\task $\begin{pmatrix} 1 & 2 \\ 4 & 5 \end{pmatrix} + \begin{pmatrix} 0 & -1 \\ -3 & 1 \end{pmatrix}$
\task $\begin{pmatrix} 1 & 0 \\ 2 & 3 \end{pmatrix} \begin{pmatrix} -1 \\ 4 \end{pmatrix}$
\task $\begin{pmatrix} 5 & -2 \\  1 & 3 \end{pmatrix} \begin{pmatrix} 0 \\ 2 \end{pmatrix}$
\end{tasks}
\end{exercisebox}
\end{frame}


%%%%%%%%%%%%%
\section[Transf. as Matrix]{Transformation using matrix}
\frame{\sectionpage}
%%%%%%%%%%%%%

\begin{frame}
\frametitle{Dilation (Scaling)}
Previously, if we only use vector notations, we have
\[ \cvec{x}{y} \to k \cvec{x}{y} = \cvec{kx}{ky},\]
which scales both coordinates by the same factor. Now, with matrix, we can scale each coordinate separately.
\[ \cvec{x}{y} \to \begin{pmatrix} \red{k} & 0 \\ 0 & \green{l} \end{pmatrix} \cvec{x}{y} = \cvec{\red{k}x}{\green{l}y}\]
$\red{k}$ and $\green{l}$ are real numbers.
\end{frame}

\begin{frame}
\frametitle{Reflection over the x-axis and y-axis}
A reflection over the x-axis is
\[ \cvec{x}{y} \to \begin{pmatrix} 1 & 0 \\ 0 & -1 \end{pmatrix} \cvec{x}{y} = \cvec{x}{-y}\]
A reflection over the y-axis is
\[ \cvec{x}{y} \to \begin{pmatrix} -1 & 0 \\ 0 & 1 \end{pmatrix} \cvec{x}{y} = \cvec{-x}{y}\]
\end{frame}

\begin{frame}
\addfig{0.9}{2dtransform_reflection}{Reflection over the x-axis (left) and over the y-axis (right)}{}
\end{frame}

\begin{frame}
\frametitle{Rotation}
\addfig{0.5}{2dtransform_rotation.png}{The rotation around the point $O$ by angle $\alpha$}{}
\end{frame}

\begin{frame}
\frametitle{Rotation}
A counterclockwise rotation around the origin by $\theta$ degree can be computed by
\[ \cvec{x}{y} \to \begin{pmatrix} \cos\theta & -\sin\theta \\ \sin\theta & \cos\theta \end{pmatrix} \cvec{x}{y}.\]
A clockwise rotation around the origin by $\theta$ degree can be computed by
\[ \cvec{x}{y} \to \begin{pmatrix} \cos(-\theta) & -\sin(-\theta) \\ \sin(-\theta) & \cos(-\theta) \end{pmatrix} \cvec{x}{y}.\]
\end{frame}

 
\begin{frame}[t] 
\begin{ex} 
Rotate the following vectors counterclockwise around the origin by $60^{\circ}$
\begin{tasks}(2)
\task $\cvec{2}{0}$ 
\task $\cvec{2\sqrt{3}}{2}$
\end{tasks}
\end{ex} 
\begin{sol}
Let $\theta = 60^{\circ}$. Then we have $\cos 60^{\circ} = \frac{1}{2}$ and $\sin 60^{\circ} = \frac{\sqrt{3}}{2}$. Hence, the transformation matrix is
\[\begin{pmatrix} \cos 60^{\circ}& -\sin 60^{\circ} \\ \sin 60^{\circ} & \cos 60^{\circ} \end{pmatrix} = \begin{pmatrix} \frac{1}{2} & -\frac{\sqrt{3}}{2} \\ \frac{\sqrt{3}}{2}& \frac{1}{2} \end{pmatrix}\]
Hence, the result of rotation is the same as the matrix-vector multiplication.
\end{sol}
\end{frame}


\begin{frame}
\begin{columns}
\begin{column}{0.5\textwidth}
\addfig{1}{2dtransform_rotation_ex1}{Rotating $\cvec{2}{0}$ by $60^{\circ}$}{}
\end{column}
\begin{column}{0.5\textwidth}
\addfig{1}{2dtransform_rotation_ex2}{Rotating $\cvec{2\sqrt{3}}{2}$ by $60^{\circ}$}{}
\end{column}
\end{columns}
\end{frame}

\begin{frame}[t]
\begin{sol} (continued)
\begin{eq*} 
\begin{pmatrix} \red{\frac{1}{2}} & \green{-\frac{\sqrt{3}}{2}} \\ \red{\frac{\sqrt{3}}{2}}& \green{\frac{1}{2}} \end{pmatrix} \cvec{2}{0} &=&2 \cvec{\red{\frac{1}{2}} }{\red{\frac{\sqrt{3}}{2}}} + 0 \cvec{\green{-\frac{\sqrt{3}}{2}}}{\green{\frac{1}{2}}} \\
&=& \cvec{2 \times \frac{1}{2}}{2 \times \frac{\sqrt{3}}{2}} + \cvec{0 \times -\frac{\sqrt{3}}{2}}{0 \times \frac{1}{2}}  \\
&=& \cvec{1}{\sqrt{3}}
\end{eq*}
\end{sol} 
\end{frame}

\begin{frame}[t]
\begin{sol} (continued)
\begin{eq*} 
\begin{pmatrix} \red{\frac{1}{2}} & \green{\frac{\sqrt{3}}{2}} \\ \red{-\frac{\sqrt{3}}{2}}& \green{\frac{1}{2}} \end{pmatrix} \cvec{2\sqrt{3}}{2} &=&2\sqrt{3} \cvec{\red{\frac{1}{2}} }{\red{\frac{\sqrt{3}}{2}}} + 2 \cvec{\green{-\frac{\sqrt{3}}{2}}}{\green{\frac{1}{2}}} \\
&=& \cvec{2\sqrt{3} \times \frac{1}{2}}{2\sqrt{3} \times \frac{\sqrt{3}}{2} } + \cvec{2\times -\frac{\sqrt{3}}{2}}{2 \times \frac{1}{2}}  \\
&=& \cvec{\sqrt{3}}{3} + \cvec{-\sqrt{3}}{1} \\
&=& \cvec{\sqrt{3}-\sqrt{3}}{3+1} = \cvec{0}{4}
\end{eq*}
\end{sol} 
\end{frame}


\begin{frame}
\frametitle{Shearing}
\addfig{0.8}{2dtransform_shear}{Shearing from the blue figure to green figure.}{}
\end{frame}

\begin{frame}
\frametitle{Shearing}
\begin{columns}
\begin{column}{0.5\textwidth}
\[ \cvec{x}{y} \to \begin{pmatrix} 1 & \red{k} \\ 0 & 1 \end{pmatrix} \cvec{x}{y} = \cvec{x+\red{k}y}{y},\]
\addfig{1}{2dtransform_shear_xaxis}{Shearing parallel to the x-axis }{}
\end{column}
\begin{column}{0.5\textwidth}
\[ \cvec{x}{y} \to \begin{pmatrix} 1 & 0 \\ \green{k} & 1 \end{pmatrix} \cvec{x}{y} = \cvec{x}{\green{k}x+y}\]
\addfig{1}{2dtransform_shear_yaxis}{Shearing parallel to the y-axis}{}
\end{column}
\end{columns}
\end{frame}

\begin{frame}
\frametitle{General Cases}
In general, given a matrix $\begin{pmatrix} a & b \\ c & d \end{pmatrix}$, the transformation
\[ \cvec{x}{y} \to \begin{pmatrix} a & b \\ c & d \end{pmatrix} \cvec{x}{y} \]
is a linear transformation. Some of them are not common, and thus do not have a name.
\end{frame}

\begin{frame}
\begin{exercisebox}{}{}
Rotate the following vectors.
\begin{tasks}(2)
\task $\begin{pmatrix} 2 \\ 0 \end{pmatrix}$ counter-clockwise by $45^\circ$.
\task $\begin{pmatrix} 3 \\ 3 \end{pmatrix}$ counter-clockwise by $45^\circ$.
\task $\begin{pmatrix} 2 \\ -1 \end{pmatrix}$ clockwise by $180^\circ$.
\task $\begin{pmatrix} 4 \\ -5 \end{pmatrix}$ clockwise by $90^\circ$.
\end{tasks}
\end{exercisebox}
\end{frame}


%%%%%%%%%%%%%
\section[Composition]{Composition of Transformations}
\frame{\sectionpage}
%%%%%%%%%%%%%


\begin{frame}
\frametitle{Motivation}
Recall that a composition of dilations in one dimension can be expressed as a single dilation. For example, we have
\[ (2)(3) x = 6x\]
Likewise, a composition of any number of linear transformations can be expressed as a single transformation. \href{https://youtu.be/XkY2DOUCWMU?t=269}{Click to see animation.}
\end{frame}

\begin{frame}
\addfig{0.9}{2dtransform_composition_intro1}{A transformation by matrix $M_1$...}{}
\end{frame}

\begin{frame}
\addfig{0.9}{2dtransform_composition_intro2}{... followed by a transformation by matrix $M_2$}{}
\end{frame}

\begin{frame}
\addfig{0.9}{2dtransform_composition_intro3}{Is that the same as a single transformation of some matrix?}{}
\end{frame}

\begin{frame}
\frametitle{Matrix Multiplication}
\begin{defn}{Matrix Multiplication}{}
\begin{equation}
\begin{pmatrix} \red{a} & \red{b} \\ \green{c} & \green{d} \end{pmatrix} \begin{pmatrix} \blue{a'} & \orange{b'} \\ \blue{c'} & \orange{d' }\end{pmatrix}= \begin{pmatrix} \red{a}\blue{a'} + \red{b}\blue{c'} & \red{a}\orange{b'} +\red{b}\orange{d'} \\ \green{c}\blue{a'} + \green{d}\blue{c'} & \green{c}\orange{b'} + \green{d}\orange{d' }\end{pmatrix}
\end{equation}
Another way to look at the computation is that the first (left) column is the product $\begin{pmatrix} \red{a} & \red{b} \\ \green{c} & \green{d} \end{pmatrix} \begin{pmatrix} \blue{a'} \\ \blue{c'}\end{pmatrix}$ and the second (right) column is $\begin{pmatrix} \red{a} & \red{b} \\ \green{c} & \green{d} \end{pmatrix} \begin{pmatrix}  \orange{b'} \\  \orange{d' }\end{pmatrix}$.
\end{defn}
\end{frame}


\begin{frame}[t]
\begin{ex}
Compute $\begin{pmatrix} 1 & 2 \\ 3 & 4 \end{pmatrix} \begin{pmatrix} 0 & 1 \\ -2 & 3\end{pmatrix}$
\end{ex}
\begin{sol}
\begin{eq*}
\begin{pmatrix} \red{1} & \red{2} \\ \green{3} & \green{4} \end{pmatrix} 
\begin{pmatrix} \blue{0} & \orange{1} \\ \blue{-2} & \orange{3}\end{pmatrix}
&=& \begin{pmatrix} \red{1}\times \blue{0} + \red{2} \times \blue{-2} & \red{1}\times \orange{1} +\red{2} \times \orange{3} \\ \green{3} \times \blue{0} + \green{4} \times \blue{-2} & \green{3} \times \orange{1} + \green{4} \times \orange{3 }\end{pmatrix} \\
&=& \begin{pmatrix} -4 & -3 \\ 7 & 15 \end{pmatrix}
\end{eq*}
\end{sol}
\end{frame}

\begin{frame}
Let $A$ and $B$ be matrices. Imagine transforming a vector $\vec{v}$ by the matrix $B$ to get $B \vec{v}$, then transforming the result by the matrix $A$ to get $A (B \vec{v})$. Instead of performing two transformations, we can use the identity
\[A(B \vec{v}) = (AB) \vec{v}\]

\frametitle{Composition}
\addfig{0.8}{2dtransform_composition}{Composition of Transformations}
\end{frame}

 
\begin{frame}[t] 
\begin{ex} 
Describe the result of the following composition, both verbally in terms of transformations, and both mathematically as a matrix.

\begin{center}
Dilation $\begin{pmatrix} 3 & 0 \\ 0 & 3 \end{pmatrix}$ and dilation $\begin{pmatrix} 2 & 0 \\ 0 & 2 \end{pmatrix}$.
\end{center}
\end{ex} 
\begin{sol}
In terms of transformation, dilating (or scaling) twice is the same as dilating once by the product of each scaling factor. In other word, the result should be a dilation by a factor of $3 \times 2 = 6$.

As a matrix, we have
\[\begin{pmatrix} 3 & 0 \\ 0 & 3 \end{pmatrix}\begin{pmatrix} 2 & 0 \\ 0 & 2 \end{pmatrix} = \begin{pmatrix} 6 & 0 \\ 0 & 6 \end{pmatrix},\]
which is exactly a dilation by a factor of 6.
\end{sol}
\end{frame}

 
\begin{frame}[t] 
\begin{ex} 
Describe the result of the following composition, both verbally in terms of transformations, and both mathematically as a matrix.

\begin{center}
Counterclockwise rotation by $45^{\circ} \begin{pmatrix} \tfrac{\sqrt{2}}{2} & -\tfrac{\sqrt{2}}{2} \\ \tfrac{\sqrt{2}}{2} & \tfrac{\sqrt{2}}{2} \end{pmatrix}$ twice.
\end{center}
\end{ex} 
\begin{sol}
In terms of transformation, a sequence of ratotations is identical to a single rotation by the sum of angles. Thus, we expect the result to be a counterclockwise rotation by $45^{\circ} + 45^{\circ} = 90^{\circ}$

As a matrix, we have
\[\begin{pmatrix} \tfrac{\sqrt{2}}{2} & -\tfrac{\sqrt{2}}{2} \\ \tfrac{\sqrt{2}}{2} & \tfrac{\sqrt{2}}{2} \end{pmatrix}\begin{pmatrix} \tfrac{\sqrt{2}}{2} & -\tfrac{\sqrt{2}}{2} \\ \tfrac{\sqrt{2}}{2} & \tfrac{\sqrt{2}}{2} \end{pmatrix}= \begin{pmatrix} 0 & -1 \\ 1 & 0 \end{pmatrix},\]
which is exactly a counterclockwise rotation by $90^{\circ}$.
\end{sol}
\end{frame}

\begin{frame}
\begin{exercisebox}{}{}
Describe the result of the following composition, both verbally in terms of transformations, and both mathematically as a matrix.
\begin{tasks}(1)
\task Dilation $\begin{pmatrix} 4 & 0 \\ 0 & 4 \end{pmatrix}$ and rotation by $90^{\circ} \quad \begin{pmatrix} 0 & -1\\ 1& 0 \end{pmatrix}$
\task Dilation $\begin{pmatrix} 2 & 0 \\ 0 & 2 \end{pmatrix}$ and rotation by $45^{\circ} \quad \begin{pmatrix} \frac{\sqrt{2}}{2} & -\frac{\sqrt{2}}{2} \\ \frac{\sqrt{2}}{2} & \frac{\sqrt{2}}{2} \end{pmatrix}$
\task Rotation by $90^{\circ} \quad \begin{pmatrix} 0 & -1 \\ 1 & 0 \end{pmatrix}$ and rotation by $90^{\circ} \quad \begin{pmatrix} 0 & -1\\ 1& 0 \end{pmatrix}$
\end{tasks}
\end{exercisebox}
\end{frame}


%%%%%%%%%%%%%
\section{Extra}
\frame{\sectionpage}
%%%%%%%%%%%%%

\begin{frame}
\frametitle{Formula for Sine and Cosine of Angle Sums}
We can prove the formula for angle sums using matrix multiplication. Consider the identity
\[\begin{pmatrix} \cos (A+B) & -\sin (A+B) \\ \sin (A+B) & \cos (A+B) \end{pmatrix} =  \begin{pmatrix} \cos A & -\sin A \\ \sin A & \cos A \end{pmatrix} \begin{pmatrix} \cos B & -\sin B \\ \sin B & \cos B \end{pmatrix}.\]
This says rotating a vector by angle $A+B$ is the same as rotating a vector by angle $A$, then rotating the vector by angle $B$. 

\end{frame}

\begin{frame}
On the other hand, we can multiply matrices on the right hand side to get
\begin{eq*}
&&\begin{pmatrix} \cos A & -\sin A \\ \sin A & \cos A \end{pmatrix} \begin{pmatrix} \cos B & -\sin B \\ \sin B & \cos B \end{pmatrix} \\
&=& \begin{pmatrix} \cos A \cos B - \sin A \sin B& -(\sin A \cos B + \cos A \sin B) \\ \sin A \cos B + \cos A\sin B & \cos A \cos B - \sin A \sin B \end{pmatrix}
\end{eq*}
By comparing the entries, we have
\begin{block}{}
\begin{eq*}
\cos (A+B) &=& \cos A \cos B - \sin A \sin B \\
\sin (A+B) &=& \sin A \cos B + \cos A \sin B
\end{eq*}
\end{block}
which are the formulas we have memorized before.
\end{frame}